% ===========================================================================
%  Homework 01 — Probability, Fall 2026.  Due September 2, 2026.
%
%  Assigned (assignment/Prob.Homework1.2026.pdf):
%    Problem 1 (instructor's own), then Ross ch.1 exercises 8, 10, 12, 39, 41.
%    Practice only, not handed in: ch.1 exercises 4, 5, 6, 7, 17.
%
%  DIVISION OF LABOUR (see AI_INSTRUCTIONS.md §7):
%    The assistant transcribes each assigned problem statement faithfully from
%    the assignment sheet in this folder's assignment/ directory, and emits an
%    empty, marked solution region beneath it.
%    The user types the mathematics inside those regions. No assistant writes
%    inside a SOLUTION region in normal mode.
% ===========================================================================
\documentclass[11pt]{article}
\usepackage{coursemacros}

\title{Homework 01}
\author{Mikey Ferguson}
\date{\today}

\begin{document}
\maketitle

\begin{problem}{1}
  Let $E$ and $F$ be events. Using the axioms of probability, prove that
  \[
    P(E \cup F) = P(E) + P(F) - P(EF).
  \]
\end{problem}
% ===== SOLUTION 1 =====
% TODO(mferguson): your work goes here.
% ===== END SOLUTION 1 =====

\begin{problem}{8}
  If $P(E) = 0.9$ and $P(F) = 0.8$, show that $P(EF) \geqslant 0.7$. In general,
  show that
  \[
    P(EF) \geqslant P(E) + P(F) - 1.
  \]
  This is known as Bonferroni's inequality.
\end{problem}
% ===== SOLUTION 8 =====
% TODO(mferguson): your work goes here.
% ===== END SOLUTION 8 =====

\begin{problem}{10}
  Show that
  \[
    P\!\left( \bigcup_{i=1}^{n} E_i \right) \leqslant \sum_{i=1}^{n} P(E_i).
  \]
  This is known as Boole's inequality.

  \upshape\emph{Hint:} Either use Equation (1.2) and mathematical induction, or
  else show that $\bigcup_{i=1}^{n} E_i = \bigcup_{i=1}^{n} F_i$, where
  $F_1 = E_1$, $F_i = E_i \bigcap_{j=1}^{i-1} E_j^{c}$, and use property (iii)
  of a probability.
\end{problem}
% ===== SOLUTION 10 =====
% TODO(mferguson): your work goes here.
% ===== END SOLUTION 10 =====

\begin{problem}{12}
  Let $E$ and $F$ be mutually exclusive events in the sample space of an
  experiment. Suppose that the experiment is repeated until either event $E$ or
  event $F$ occurs. What does the sample space of this new super experiment look
  like? Show that the probability that event $E$ occurs before event $F$ is
  $P(E) / [P(E) + P(F)]$.

  \upshape\emph{Hint:} Argue that the probability that the original experiment is
  performed $n$ times and $E$ appears on the $n$th time is
  $P(E) \times (1-p)^{n-1}$, $n = 1, 2, \ldots$, where $p = P(E) + P(F)$. Add
  these probabilities to get the desired answer.
\end{problem}
% ===== SOLUTION 12 =====
% TODO(mferguson): your work goes here.
% ===== END SOLUTION 12 =====

\begin{problem}{39}
  Stores $A$, $B$, and $C$ have 50, 75, and 100 employees and, respectively, 50,
  60, and 70 percent of these are women. Resignations are equally likely among
  all employees, regardless of sex. One employee resigns and this is a woman.
  What is the probability that she works in store $C$?
\end{problem}
% ===== SOLUTION 39 =====
% TODO(mferguson): your work goes here.
% ===== END SOLUTION 39 =====

\begin{problem}{41}
  In a certain species of rats, black dominates over brown. Suppose that a black
  rat with two black parents has a brown sibling.
  \begin{enumerate}
    \item[(a)] What is the probability that this rat is a pure black rat (as
      opposed to being a hybrid with one black and one brown gene)?
    \item[(b)] Suppose that when the black rat is mated with a brown rat, all
      five of their offspring are black. Now, what is the probability that the
      rat is a pure black rat?
  \end{enumerate}
\end{problem}
% ===== SOLUTION 41 =====
% TODO(mferguson): your work goes here.
% ===== END SOLUTION 41 =====

\end{document}
